\section*{Bab \ref{ch:g12:catalysis} --- Katalisis}

\begin{solution}{exo:g12:catalysis:1}
\omterm{def:g7:chemical-reactions:reactant}{Pereaksi}: \ce{H2O2}. Produk: \ce{H2O} dan \ce{O2}. \omterm{def:g12:catalysis:catalyst}{Katalis}: \ce{I-}
(dipakai pada tahap 1, dikembalikan pada tahap 2). Zat antara:
\ce{IO-} (terbentuk pada tahap 1, dipakai pada tahap 2).
\end{solution}

\begin{solution}{exo:g12:catalysis:2}
Homogen; heterogen; heterogen; homogen.
\end{solution}

\begin{solution}{exo:g12:catalysis:3}
Jalan tanpa \omterm{def:g12:catalysis:catalyst}{katalis}. Tingkat awal dan tingkat akhirnya sama: produk
yang sama terbentuk, dengan energi yang dilepaskan sama; hanya jalannya
yang berbeda.
\end{solution}

\begin{solution}{exo:g12:catalysis:4}
\omterm{def:g12:catalysis:catalyst}{Katalis} dipakai pada satu tahap dan dikembalikan pada tahap lain:
\omterm{def:g12:catalysis:catalyst}{katalis} muncul di kedua ruas jumlah tahap-tahapnya, dan saling
meniadakan.
\end{solution}

\begin{solution}{exo:g12:catalysis:5}
Reaksinya terjadi di permukaan: makin luas permukaannya, makin banyak
\omterm{def:g7:atoms-and-molecules:molecule}{molekul} yang dapat bereaksi sekaligus, untuk massa \omterm{def:g12:catalysis:catalyst}{katalis} (yang sering
mahal) yang sama.
\end{solution}

\begin{solution}{exo:g12:catalysis:6}
Jumlahnya: \ce{2Fe^{3+} + 2H2O2 + 2Fe^{2+} + 2H+ -> 2Fe^{2+} + O2 + 2H+ +
2Fe^{3+} + 2H2O}. \omterm{def:g9:ions:ion}{Ion-ion} besi dan \omterm{def:g9:ions:ion}{ion-ion} hidrogen muncul di kedua
ruas dan saling meniadakan: \ce{2H2O2 -> 2H2O + O2}.
\end{solution}

\begin{solution}{exo:g12:catalysis:7}
Setengah dari \qty{72}{mL} adalah \qty{36}{mL}: tercapai setelah sekitar
\qty{35}{min} tanpa \omterm{def:g12:catalysis:catalyst}{katalis} dan \qty{3.5}{min} dengan \omterm{def:g12:catalysis:catalyst}{katalis}. \omterm{def:g12:catalysis:catalyst}{Katalis}
memperpendeknya sepuluh kali.
\end{solution}

\begin{solution}{exo:g12:catalysis:8}
Reaksinya berlangsung di permukaan platina; timbal menutupinya dan
tetap di sana, sehingga gas-gas tidak lagi dapat mencapai \omterm{def:g9:periodic-table-first-look:metal}{logamnya}.
\omterm{def:g12:catalysis:catalyst}{Katalisnya} ``teracuni'': ada, tetapi tidak berguna.
\end{solution}

\begin{solution}{exo:g12:catalysis:9}
Sampai sekitar \qty{37}{\celsius}, reaksinya makin cepat dengan naiknya
suhu, seperti reaksi apa pun. Di atas suhu itu, \omterm{def:g12:catalysis:enzyme}{enzimnya}, sebuah
protein, kehilangan bentuknya dan bersama itu kehilangan aktivitasnya;
pada \qty{80}{\celsius} \omterm{def:g12:catalysis:enzyme}{enzim} itu rusak.
\end{solution}

\begin{solution}{exo:g12:catalysis:10}
\ce{2CO + 2NO -> 2CO2 + N2}: C 2 = 2, O 4 = 4, N 2 = 2. Karbon
monoksida dioksidasi oleh nitrogen monoksida, yang tereduksi: setiap
polutan menyingkirkan polutan yang lain.
\end{solution}

\begin{solution}{exo:g12:catalysis:11}
$n(\ce{O2}) = 0.072 / 24.1 = \qty{2.99e-3}{mol}$;
$n(\ce{H2O2}) = 2 \times \num{2.99e-3} = \qty{5.98e-3}{mol}$ dalam
\qty{10.0}{mL}: $c = \qty{0.598}{mol/L}$.
\end{solution}

\begin{solution}{exo:g12:catalysis:12}
\ce{C2H5OH -> C2H4 + H2O}, sebuah \omterm{def:g12:curly-arrows:categories}{eliminasi}; \ce{C2H5OH -> CH3CHO + H2}.
\omterm{def:g12:catalysis:selective}{Katalis selektif} mempercepat salah satu dari beberapa reaksi yang
mungkin jauh lebih banyak daripada yang lain, sehingga menentukan
produknya.
\end{solution}

\begin{solution}{exo:g12:catalysis:13}
\omterm{def:g12:catalysis:catalyst}{Katalis} tidak mengubah \omterm{def:g10:reaction-progress-table:system}{keadaan akhir}: kurva dengan dan tanpa \omterm{def:g12:catalysis:catalyst}{katalis}
berakhir pada volume oksigen yang sama. Jumlah produk yang diperoleh
sama, hanya lebih cepat.
\end{solution}

\begin{solution}{exo:g12:catalysis:14}
$\num{6.02e23} / \num{5e6} = \qty{1.2e17}{s}$, sekitar empat miliar
tahun. Untuk menyelesaikannya dalam satu detik: \num{1.2e17} \omterm{def:g7:atoms-and-molecules:molecule}{molekul}
\omterm{def:g12:catalysis:enzyme}{enzim}, yaitu \qty{2e-7}{mol}.
\end{solution}

\begin{solution}{exo:g12:catalysis:15}
Amonia \ce{NH3}: $150 \times 17.0 / 14.0 = \qty{182}{Mt}$. \omterm{def:g7:atoms-and-molecules:atom}{Atom}
nitrogen: $\num{1.50e14} / 14.0 = \qty{1.07e13}{mol}$, yaitu
\qty{5.36e12}{mol} \ce{N2}.
\end{solution}

\begin{solution}{pb:g12:catalysis:1}
\textbf{1.} \ce{2H2O2 -> 2H2O + O2}.

\textbf{2.} Ya: hidrogen peroksida adalah \omterm{def:g11:redox:oxidant}{oksidator} pasangan
\ce{H2O2}/\ce{H2O} dan \omterm{def:g11:redox:oxidant}{reduktor} pasangan \ce{O2}/\ce{H2O2}; zat itu
mengoksidasi dan mereduksi dirinya sendiri.

\textbf{3.} Tanpa \omterm{def:g12:catalysis:catalyst}{katalis}, reaksinya sangat lambat.

\textbf{4.} $20 / 22.4 = \qty{0.893}{mol}$.

\textbf{5.} Dua \omterm{def:g10:the-mole:amount}{mol} hidrogen peroksida per \omterm{def:g10:the-mole:amount}{mol} oksigen:
\qty{1.79}{mol}.

\textbf{6.} \qty{1.79}{mol/L}.

\textbf{7.} $M = 2 \times 1.0 + 2 \times 16.0 = \qty{34.0}{g/mol}$;
$C_m = 1.79 \times 34.0 = \qty{60.7}{g/L}$.

\textbf{8.} Setengahnya: \qty{0.893}{mol/L}.

\textbf{9.} Katalase: sebuah \omterm{def:g12:catalysis:enzyme}{enzim}, terlarut, \omterm{def:g12:catalysis:catalyst}{katalis} biologis; \omterm{def:g9:ions:ion}{ion}
besi(III): \omterm{def:g12:catalysis:homogeneous}{katalisis homogen}.

\textbf{10.} \ce{2Fe^{3+} + H2O2 -> 2Fe^{2+} + O2 + 2H+} dan
\ce{2Fe^{2+} + H2O2 + 2H+ -> 2Fe^{3+} + 2H2O}; jumlahnya adalah
\ce{2H2O2 -> 2H2O + O2}.

\textbf{11.} Dengan \omterm{def:g12:catalysis:catalyst}{katalis}, volumenya naik jauh lebih cepat; kedua
kurva mendatar pada volume akhir yang sama.

\textbf{12.} \omterm{def:g9:ions:ion}{Ion-ion} besi(III) (jingga) diubah menjadi besi(II) (hijau
pucat) pada tahap pertama dan kembali menjadi besi(III) pada tahap
kedua; pada akhirnya, seluruh besi kembali menjadi besi(III).

\textbf{13.} Perebusan telah merusak \omterm{def:g12:catalysis:enzyme}{enzimnya}.

\textbf{14.} $0.100 \times 20 = \qty{2.0}{L}$.

\textbf{15.} $n(\ce{O2}) = 0.100 \times 0.893 = \qty{0.0893}{mol}$,
yaitu $0.0893 \times 24.1 = \qty{2.15}{L}$.

\textbf{16.} Penguraian yang lambat melepaskan oksigen, yang akan
menimbulkan tekanan di dalam botol yang rapat.

\textbf{17.} $1.50 / 2 = \qty{0.750}{mol}$ oksigen per liter, yaitu
$0.750 \times 22.4 = \qty{16.8}{L}$: ``16.8 volume''.

\textbf{18.} $(1.79 - 1.50)/1.79 \approx \qty{16}{\percent}$.

\textbf{19.} \qty{1.79}{mol/L}.
\end{solution}
